\textbf{Гоголь} (падает из-за кулис на сцену и смирно лежит).

\textbf{Пушкин} (выходит, спотыкается об Гоголя и падает): Вот черт! Никак об Гоголя!

\textbf{Гоголь} (поднимаясь): Мерзопакость какая! Отдохнуть не дадут! (Идет, спотыкается об Пушкина и падает). Никак об Пушкина спотыкнулся!

\textbf{Пушкин} (поднимаясь): Ни минуты покоя! (Идет, спотыкается об Гоголя и падает). Вот черт! Никак опять об Гоголя!

\textbf{Гоголь} (поднимаясь): Вечно во всем помеха! (Идет, спотыкается об Пушкина и падает). Вот мерзопакость! Опять об Пушкина!

\textbf{Пушкин} (поднимаясь): Хулиганство! Сплошное хулиганство! (Идет, спотыкается об Гоголя и падает). Вот черт! Опять об Гоголя!

\textbf{Гоголь} (поднимаясь): Это издевательство сплошное! (Идет, спотыкается об Пушкина и падает). Опять об Пушкина!

\textbf{Пушкин} (поднимаясь): Вот черт! Истинно что черт! (Идет, спотыкается об Гоголя и падает). Об Гоголя!

\textbf{Гоголь} (поднимаясь): Мерзопакость! (Идет, спотыкается об Пушкина и падает). Об Пушкина!

\textbf{Пушкин} (поднимаясь): Вот черт! (Идет, спотыкается об Гоголя и падает за кулисы). Об Гоголя!

\textbf{Гоголь} (поднимаясь): Мерзопакость! (Уходит за кулисы).

За сценой слышен голос Гоголя: <<Об Пушкина!>>

Занавес.

1934 г.